\chapter{Teichm\"uller space} \label{Teichmuller:chapter}
We have discovered in Chapter \ref{surfaces:chapter} that every surface $S_g$ of genus $g\geqslant 2$ can be equipped with a hyperbolic metric, and we have already noticed that this metric is not unique: this chapter is entirely devoted to studying this non-uniqueness phenomenon.

It turns out that $S_g$ admits a continuous family of non-equivalent hyperbolic metrics, that form altogether a nice topological space called the \emph{Teichm\"uller space} of $S_g$. We prove in his chapter that the Teichm\"uller space of $S_g$ is homeomorphic to an open ball of dimension $6g-6$. To prove this fact we will introduce and study concepts like \emph{length functions}, \emph{earthquakes}, and \emph{Fenchel--Nielsen coordinates}. The simple closed curves and their geodesic representatives will play a fundamental role in all the discussion.

\section{Introduction}
Let $S_g$ be as usual a closed orientable surface of genus $g$. We know that $S_g$ admits an elliptic, flat, or hyperbolic metric if and only if $g=0$, $g=1$, or $g\geqslant 2$ respectively. The elliptic metric on the two-sphere is unique up to isometries, but the flat and hyperbolic metrics on the other surfaces are not.

We want to define the space of all flat or hyperbolic metrics on $S_g$ when $g\geqslant 1$. There are two natural ways do to this:\index{Teichm\"uller space}\index{moduli space}

\begin{defn}
The \emph{moduli space} of $S_{g}$ is the set of all the flat or hyperbolic metrics on $S_{g}$ considered up to orientation-preserving isometries and rescaling.

The \emph{Teichm\"uller space} $\Teich(S_g)$ is the set of all the flat or hyperbolic metrics on $S_{g}$ considered up to isometries isotopic to the identity and rescaling.
\end{defn}

The \emph{rescaling} of the metric is a simple operation that takes place only on the flat metrics on the torus $T$. On the torus $T$ a flat metric $g$ can be rescaled by any constant $\lambda >0$ to give another flat metric $\lambda g$: the rescaling changes the lengths by a factor $\sqrt\lambda$ and the area by a factor $\lambda$. Up to rescaling, we may for instance require that $T$ has unit area.

At a first sight, the moduli space seems a more natural object to study. It turns out however that the Teichm\"uller space is homeomorphic (for some natural topology) to an open ball, while the moduli space is topologically more involved: it is then more comfortable to define and study the Teichm\"uller space first, and then consider the moduli space as a quotient of Teichm\"uller space.

\subsection{The Teichm\"uller space of the torus}
The flat metrics on the torus $T$ are classified quite easily. We have seen in Proposition \ref{flat:is:torus:prop} that every flat torus $T$ is of type $\matC/_\Gamma$ for some lattice $\Gamma <\matC$ isomorphic to $\matZ^2$. 

Fix two generators $m, l$ for $\pi_1(T)$. These are identified to two generators of $\Gamma$. Up to rescaling the metric, rotating $\matC$ around the origin, and reflecting along the real axis, we may suppose that these generators are 1 and some complex number $z$ lying in the upper half-plane $H^2$, so that $\Gamma = \langle 1,z\rangle$.

\begin{prop} \label{Teich:T:prop}
By sending the flat metric on $T$ to $z$ we get a bijection:
$$ \Teich(T) \longrightarrow H^2.$$
\end{prop}
\begin{proof}
The map is well-defined: two metrics related by an isometry isotopic to the identity produce the same $z$. The inverse $H^2 \to \Teich(T)$ is constructed by identifying $T$ with $\matC/_{\langle 1,z\rangle}$ sending $(m,l)$ to $(1,z)$.
\end{proof}

\begin{figure}
\begin{center}
\includegraphics[width = 8cm] {\iftoggle{BW}{toro_dominio-BW}{toro_dominio}}
\nota{The flat metric on the torus $T$ determined by $z\in H^2$ may be constructed by identifying the opposite sides of the parallelogram with vertices $0,1,z,z+1$. The lattice $\Gamma$ is generated by $1$ and $z$ and the parallelogram is a fundamental domain.}
\label{toro_dominio:fig}
\end{center}
\end{figure}

The flat metric that corresponds to $z \in H^2$ may be constructed by identifying the opposite sides of a parallelogram as in Figure \ref{toro_dominio:fig}. We will often tacitly identify $\Teich(T)$ with the projective plane $H^2$ via this bijective correspondence.

\begin{oss} \label{torus:isometries:oss}
Let $T=\matC/_\Gamma$ be a flat torus. Every translation $z \mapsto z +w$ in $\matC$ commutes with $\Gamma$ and hence descends to an isometry of $T$. Therefore the isometry group $\Iso^+(T)$ is not discrete: every flat torus is \emph{homogeneous}, \emph{i.e.}~for every pair of points $x,y \in T$ there is an isometry sending $x$ to $y$. 
\end{oss}

\begin{oss}
Let $T=\matC/_\Gamma$ be a flat torus. Every non-trivial element $\gamma\in\pi_1(T)$ is represented by a closed geodesic, unique up to translations. Its counterimage in $\matC$ consists of parallel lines whose slope depends only on $\gamma$.  The geodesic is simple if and only if $\gamma$ is primitive. 
\end{oss}

\subsection{Action of the mapping class group}
Recall that the mapping class group $\MCG(S_g)$ of $S_g$ is the group of all the orientation-preserving self-diffeomorphisms of $S_g$ considered up to isotopy (or equivalently, homotopy). We now show that the mapping class group of $S_g$ acts on its Teichm\"uller space.

A diffeomorphism $\varphi\colon S_g \to S_g$ transports a metric $m$ on $S_g$ into a new metric $\varphi_*m$ by pushing it forward as follows:
$$(\varphi_*m)_{\varphi(x)}\big(d\varphi_x(v), d\varphi_x(w)\big) = m_x(v,w).$$
If $m$ varies through an isotopy, the metric $\varphi_*m$ varies through a corresponding isotopy: therefore $\varphi$ acts on $\Teich(S_g)$ as follows
\begin{align*}
\Teich(S_g) & \longrightarrow \Teich(S_g) \\
[m] & \longmapsto [\varphi_*m]
\end{align*} 
If we vary $\varphi$ by an isotopy the action is unaffected. 
Therefore the mapping class group $\MCG(S_g)$ acts on $\Teich(S_g)$, and by definition the quotient
$$\Teich(S_g)/_{\MCG(S_g)}$$
is the moduli space of $S_g$. 

Everything can be written explicitly for the torus. Recall that we have identified $\Teich(T)$ with $H^2$ and $\MCG(T)$ with $\SLZ$, see Proposition \ref{Torelli:toro:prop}.

\begin{prop}  \label{toro:azione:prop}
The action of $\MCG(T)$ on $\Teich(T)$ is the following action of $\SLZ$ on $H^2$ as M\"obius transformations: 
$$\begin{pmatrix} a & b \\ c & d \end{pmatrix} \colon z \longmapsto \frac{az-b}{-cz+d}.$$
\end{prop}
\begin{proof}
The metric $z$ assigns to $T$ the structure $\matR^2/_\Gamma$ with $\Gamma = \langle 1,z \rangle$ and $(m,l)$ sent to $(1,z)$. 

Pick $\varphi =  \left(\begin{smallmatrix} a & b \\ c & d \end{smallmatrix}\right) \in \SLZ =\MCG(T)$. Since $\varphi^{-1} = \left(\begin{smallmatrix} d & -b \\ -c & a \end{smallmatrix}\right)$, in the new metric $\varphi_*$ we assign to $(m,l)$ the translations $(d-cz, -b+az)$, which transform via rotations and dilations into $(1,  \frac{az-b}{-cz+d}).$
\end{proof}

We note in particular that $\MCG(T)$ acts via isometries on the hyperbolic plane $H^2$. The kernel of the action is $\{\pm I\}$: two matrices $A$ and $-A$ act in the same way on $\Teich(T)$. 

\begin{cor}
The moduli space of $T$ is the orbifold $H^2/_{\PSLZ}$.
\end{cor}

The orbifold $H^2/_{\PSLZ}$ is described in Figure \ref{orbifold_Farey:fig}. It has two singular points of order 2 and 3: these represent the square and hexagonal torus, see Section \ref{piatte:subsection} and Figure \ref{modular_group:fig}.

We have constructed a bijection between $\Teich(T)$ and $H^2$, and we now want to construct for $g\geqslant 2$ an analogous identification between $\Teich(S_g)$ and some open set of $\matR^N$ for some $N$ depending on $g$. To this purpose we need to introduce some concepts.

\section{Earthquakes and length functions}
Simple closed geodesics are a formidable tool to study $\Teich(S_g)$: we can use a simple closed geodesic to \emph{twist} a metric (the operation is called an \emph{earthquake}), and by simply looking at the lengths of the other closed geodesics we can measure how the metric varies along this transformation. We introduce these operations here; later on, we will use them to parametrize $\Teich(S_g)$.

\subsection{Earthquakes} \label{earthquakes:subsection}
The hyperbolic, flat, and spherical metrics on surfaces may be twisted along simple closed geodesics: this operation is called an \emph{earthquake}.\index{earthquake} 

Let $m$ be a complete hyperbolic, flat, or elliptic metric on an oriented surface $S$ and $\gamma$ be a simple closed geodesic in $S$. Recall that $\gamma$ is a map $\gamma \colon S^1 \to S$. Fix an angle $\theta \in \matR$. Informally, a new complete hyperbolic, flat, or elliptic metric $m_\theta$ on $S$ is constructed by cutting $S$ along $\gamma$ and regluing with a counterclockwise twist of angle $\theta$. Formally, the new metric is defined as follows.

\begin{figure}
\begin{center}
\includegraphics[width = 10cm] {\iftoggle{BW}{prodotto-BW}{prodotto}}
\nota{A $R$-annulus around a simple closed geodesic $\gamma$ on a hyperbolic surface is the quotient of a $R$-neighbourhood of a line $l$ by a hyperbolic transformation. The orthogonal (\iftoggle{BW}{light grey}{green}) geodesic segments are parametrized by arc length as $[-R,R]$, hence the $R$-annulus is naturally parametrized as $S^1 \times [-R,R]$.}
\label{prodotto:fig}
\end{center}
\end{figure}

In the hyperbolic case, recall from Proposition \ref{tube:prop} that $\gamma$ has a $R$-neighbourhood isometric to a $R$-tube for some $R>0$. A $R$-tube in dimension two is a $R$-annulus as in Figure \ref{prodotto:fig}, defined by quotienting a $R$-neighbourhood of a line $l$ by a hyperbolic transformation. The $R$-annulus is naturally parametrized as $S^1\times [-R,R]$, where $\{e^{it}\}\times [-R,R]$ is the geodesic segment orthogonal to $\gamma$ in $\gamma(e^{it})$ parametrized by arc length.
The flat and elliptic cases are analogous.

\begin{figure}
\begin{center}
\includegraphics[width = 10cm] {\iftoggle{BW}{twist-BW}{twist}}
\nota{To define the earthquake we pick a diffeomorphism of the $R$-annulus that modifies the orthogonal segments as shown here.}
\label{twist:fig}
\end{center}
\end{figure}

We choose a diffeomorphism $\varphi$ of $S^1\times [-R, R]$ that curves the segments left-wise with step $\theta$ as in Figure \ref{twist:fig}-(right).  More precisely, let $f\colon[-R, R] \to \matR$ be a smooth function which is zero on $\left[-R, -\frac R 2\right]$ and is constantly $\theta$ on $\left[\frac R 2,R\right]$. We set $\varphi(e^{it},s) = (e^{i(t+f(s))},s)$.

We define a new metric $m_\theta$ on $S_g$ as follows: the metric tensor $m_\theta$ coincides with $\varphi_* m$ on the $R$-annulus and coincides with $m$ on the complement of the $\frac{R}2$-annulus $S^1\times \left[-\frac R2, \frac R2\right]$.

\begin{prop} The metric tensor $m_\theta$ is well-defined and gives a complete hyperbolic, flat, or elliptic metric to $S_g$. 
\end{prop}
\begin{proof}
It is well-defined because $m$ and $m_\theta$ coincide on $S^1\times \left[\frac R 2, R\right]$, since $(e^{it},s)\mapsto (e^{i(t+\theta)},s)$ is an isometry of the $R$-annulus. It is hyperbolic, flat, or elliptic because both patches $m$ and $\varphi_*m$ are.
\end{proof}

\begin{oss} \label{R:oss}
In the new metric $m_\theta$ the curve $\gamma$ is still a simple closed geodesic of the same length as before, and its $R$-neighbourhood is also unchanged. 
\end{oss}

Of course by deforming an elliptic metric in this way we get nothing new, because all the elliptic metrics on $S^2$ are isometric. Earthquakes are interesting only in the flat and hyperbolic geometries.

\subsection{The earthquake map}
We consider the surface $S_g$ with $g\geqslant 1$ and show that earthquakes define nice actions on the Teichm\"uller space $\Teich(S_g)$. 

If $m$ is any flat or hyperbolic metric and $\gamma$ is a non-trivial simple closed curve in $S_g$, we define $m_\theta^\gamma$ to be the flat or hyperbolic metric obtained from $m$ via an earthquake of angle $\theta$ performed along the unique (up to translations if $g=1$) simple closed geodesic homotopic to $\gamma$ in the metric $m$.\index{earthquake map}

\begin{prop}
The \emph{earthquake map}
\begin{align*}
E_\gamma\colon \matR \times \Teich(S_g) & \longrightarrow \Teich(S_g) \\
 (\theta, m) & \longmapsto m_\theta^{\gamma} 
\end{align*}
is a well-defined action of $\matR$ on $\Teich(S_g)$. The map $E_\gamma$ depends only on the homotopy class of $\gamma$.
\end{prop}
\begin{proof}
The only ambiguity in the definition of $m_\theta^\gamma$ is the choice of the function $f\colon [-R,R] \to \matR$. If we use another function $f'$ the resulting metric $(m_\theta^\gamma)'$ changes only by an isotopy: the diffeomorphism of $S_g$ which is the identity outside the $R$-annulus and sends $(e^{it},s)$ to $(e^{i(t+f(s)-f'(s))},s)$ is an isometry between $m_\theta^\gamma$ and $(m_\theta^\gamma)'$, and is clearly isotopic to the identity.

To prove that $E_\gamma$ is an action we need to check that 
$$m_{\theta+\theta'}^{\gamma} = \left(m_{\theta'}^\gamma\right)_\theta^\gamma.$$
By Remark \ref{R:oss} we can take the same $R$-annulus to compose two earthquakes and the equality follows.
\end{proof}

Like the Dehn twists $T_\gamma$ defined in Section \ref{Dehn:twist:subsection}, the action $E_\gamma$ depends on the orientation of $S_g$ but not on the orientation of $\gamma$. There is indeed a strong relation between earthquakes and Dehn twists on $\gamma$: as objects acting on Theichm\"uller space, the first generalise the second.

\begin{prop} \label{root:prop}
We have $T_\gamma(m) = E_\gamma(2\pi,m)$.
\end{prop}
\begin{proof}
It follows directly from the definitions.
\end{proof}

We now want to study the Teichm\"uller space and the action of $E_\gamma$ on it. In mathematics a space is often beautifully described by some natural functions defined on it: this role is played here by the length functions of closed curves.

\subsection{Length functions}
A homotopically nontrivial (possibly non simple) closed curve $\gamma$ in $S_g$ defines a \emph{length function}\index{length function}
$$\ell^\gamma\colon \Teich(S_g) \to \matR_{>0}$$
which assigns to a metric $m\in\Teich(S_g)$ the length $\ell^\gamma(m)$ of the unique closed geodesic homotopic to $\gamma$. 

When $g=1$ we must actually specify a couple of things in the definition: the closed geodesic $\gamma$ is unique only up to translations, which do not affect its length; 
on the other hand rescaling \emph{does} affect lengths, so to get a well-defined length function we always rescale the metric $m$ to have unit area. 

We want to study these length functions, and as usual we start by analysing the simpler torus flat world, where everything can be described explicitly.

\subsection{Length functions on the torus} \label{length:torus:subsection}
We denote every (isotopy class of) non-trivial simple closed curve on $T$ with a coprime pair $(p,q)$ of integers and we identify the Teichm\"uller space $\Teich(T)$ with $H^2\subset\matC$, see Propositions \ref{mn:prop} and \ref{Teich:T:prop}. The length functions may be written explicitly.

\begin{prop} \label{lunghezza:toro:prop}
The formula holds:
$$\ell^{(p,q)}(z) = \frac{|p+qz|}{\sqrt{\Im z }}$$
for every simple closed curve $(p,q)$ and every metric $z\in H^2$.
\end{prop}
\begin{proof}
Up to rescaling we have $T=\matR^2/_{\Gamma}$ with $\Gamma = \langle 1,z \rangle $. The translation in $\Gamma$ corresponding to $(p,q)$ is $p\cdot 1 + q\cdot z$ and the closed geodesic it produces has length $|p+qz|$. The area of the torus $T$ is $\Im z$ (see a fundamental domain in Figure \ref{toro_dominio:fig}) and hence we must rescale it by $1/\sqrt{\Im z}$.
\end{proof}

\begin{figure}
\begin{center}
\includegraphics[width = 9cm] {\iftoggle{BW}{toro_terremoto-BW}{toro_terremoto}}
\nota{A torus with metric $z$ (left) twisted along the horizontal curve $\gamma$ (right). The curve $\alpha$ is a closed geodesic in the new metric.}
\label{toro_terremoto:fig}
\end{center}
\end{figure}

We can also write the earthquake action on the meridian $m=(1,0)$ of $T$.

\begin{ex} \label{terremoto:toro:ex}
We have:
$$E_m(\theta, z) = z+\frac {\theta}{2\pi}.$$
\end{ex}
\begin{proof}[Hint]
Look at the closed geodesic $\alpha$ in Figure \ref{toro_terremoto:fig}.
\end{proof}

We have discovered in particular that $E_m$ is a parabolic isometry of $H^2$, and we can now consider the earthquake action along a generic curve $(p,q)$.

\begin{cor} 
The earthquake action $E_{(p,q)}$ is the 1-parameter family of parabolic transformations with fixed point $-\frac pq \in \partial H^2$. 
\end{cor}
\begin{proof}
We know this when $(p,q)=(1,0)=m$. In general, we send $(1,0)$ to $(p,q)$ via some element of the mapping class group: this element acts by isometries of $H^2$ by Proposition \ref{toro:azione:prop} and sends $\infty$ to $-\frac pq$. 
\end{proof}

The orbits of $E_{(p,q)}$ are the horospheres centred at $-\frac pq$. 
We have discovered that, quite unexpectedly, the hyperbolic geometry of the plane $H^2$ is well designed to model the Teichm\"uller space of the flat torus $T$. 

It is now natural to identify the \emph{unoriented} simple closed curve $\pm (p,q)$ with the rational point $-\frac pq$ in $\matR\cup\{\infty\} \subset \partial H^2$. We get the following convexity property.

\begin{cor} \label{convex:cor}
Let $\gamma$ and $\eta$ be two simple closed curves on $T$.
If $i(\gamma, \eta)>0$ then the length function $\ell^\eta$ is strictly convex on the orbits of $E_{\gamma}$.
\end{cor}
\begin{proof}
We may suppose $\gamma=m=(1,0)$ and note that the condition $i(\gamma,\eta)>0$ translates into $\eta = (p,q) \neq (\pm 1,0)$. The function in Proposition \ref{lunghezza:toro:prop} is strictly convex on the horospheres $\Im z = k$ when $(p,q)\neq (\pm 1,0)$.
\end{proof}

Summing up, the torus picture is the following: 
\begin{itemize}
\item the mapping class group acts on the Teichm\"uller space roughly like $\PSLZ$ acts isometrically on the hyperbolic plane $H^2$, 
\item the unoriented simple closed curves are in 1-1 correspondence with the rational points in $\partial H^2 = \matR\cup \infty$, 
\item earthquakes act like parabolic transformations centred at these rational points, \item the length function $\ell^\gamma$ is constant at the horospheres centred at $\gamma$ but strictly convex at the horospheres centred at all the other curves.
\end{itemize}

The main goal of Chapters \ref{Teichmuller:chapter} and \ref{automorfismi:chapter} is to draw a similar picture for surfaces $S_g$ of higher genus $g\geqslant 2$. Everything is more difficult in the hyperbolic world, because there are no nice explicit formulas describing the length functions, the action of the mapping class group, the simple closed curves, and the earthquakes. We now start by generalising the last point: the strict convexity of length functions.

\subsection{Convexity of the length functions} \label{convexity:lengths:subsection}
We consider $S_g$ with $g\geqslant 2$. Our aim now is to generalise Corollary \ref{convex:cor} to the hyperbolic setting and then later use this convexity property to parametrize $\Teich(S_g)$. We need a preliminary result.\index{convexity of the length functions}

\begin{ex} \label{convessa:ex}
Let $f\colon \matR^m \times \matR^n \to \matR_{\geqslant 0}$ be strictly convex and proper. The function
\begin{align*}
F\colon \matR^n & \longrightarrow \matR_{\geqslant 0} \\
y & \longmapsto \min \left\{f(x,y)\ \big| \ x \in \matR^m \right\}
\end{align*}
is well-defined, strictly convex, and proper. 
\end{ex}

The following generalisation of Corollary \ref{convex:cor} says that length functions are convex on the earthquakes orbits, and very often they are strictly convex and proper.

\begin{prop} \label{strettamente:convessa:prop}
Let $\eta$ and $\gamma$ be two homotopically non-trivial simple closed curves in $S_g$ and $m$ be a hyperbolic metric on $S_g$. The function 
\begin{align*}
\matR & \longrightarrow \matR_{\geqslant 0} \\
\theta & \longmapsto \ell^\eta(m_\theta^\gamma)
\end{align*}
is
\begin{itemize}
\item constant if $i(\eta,\gamma)=0$,
\item strictly convex and proper if $i(\eta, \gamma)>0$.
\end{itemize}
\end{prop}
\begin{proof}
The metric $m_\theta^\gamma$ is obtained by twisting $m$ of an angle $\theta$ along the simple closed geodesic $\gamma$. If $i(\eta,\gamma)=0$ the curves $\eta$ and $\gamma$ are disjoint geodesics in $m$ and $\eta$ is not affected by the earthquakes that we perform near $\gamma$, hence the function $\ell^\eta(m_\theta^\gamma)$ is constant.

Consider the case $n=i(\eta,\gamma)>0$. Denote by $\bar{\eta}^\theta$ the geodesic representative of $\eta$ in the twisted metric $m_\theta^\gamma$: it intersects $\gamma$ transversely in $n$ points. 

\begin{figure}
\begin{center}
\includegraphics[width = 10cm] {\iftoggle{BW}{spezza_new-BW}{spezza_new}}
\nota{A geodesic in $m_\theta^\gamma$ can be seen on the original metric $m$ as follows: it is a geodesic outside the annulus and deviates on the left by an angle $\theta$ every time it crosses it (centre). We may simplify the picture by describing it as a broken geodesic line that makes a left $\theta$-jump each time that it crosses $\gamma$ (right).}
\label{spezza:fig}
\end{center}
\end{figure}

Fix a sufficiently small $R$-annulus around $\gamma$ and note that the geodesics in $m_\theta^\gamma$ can be seen in the original metric $m$ as follows: these are curves that are geodesic outside the $R$-annulus and deviate smoothly on the left each time they cross it as in Figure \ref{spezza:fig}-(centre).
We may substitute each smooth deviation with a broken jump as shown in Figure \ref{spezza:fig}-(right) and get a bijection
$$\left\{\begin{array}{c} {\rm closed\ geodesics} \\ {\rm with\ respect\ to\ } m_\theta^\gamma \end{array} \right\}
\longleftrightarrow \left\{\begin{array}{c} {\rm broken\ geodesics} \\ {\rm with\ respect\ to\ } m \end{array}\right\}$$
where a \emph{broken geodesic} is a geodesic that at every crossing of $\gamma$ jumps to the left at distance $\frac{\theta L(\gamma)}{2\pi}$ and then keeps going on, leaving $\gamma$ with the same incidence angle (here $L(\gamma)$ is the length of $\gamma$). This correspondence is useful because it preserves the lengths: the length of the closed geodesic for $m_\theta^\gamma$ is equal to the length of the corresponding broken geodesic (which is the sum of the lengths of its components), because the segments in Figure \ref{spezza:fig}-(left) and (right) are isometric.

\begin{figure}
\begin{center}
\includegraphics[width = 6cm] {\iftoggle{BW}{spezza_su-BW}{spezza_su}}
\nota{The \iftoggle{BW}{}{(blue) }line $l$ is a lift of $\bar \eta$. Consider $n$ consecutive intersections with lifts $r_1,\ldots, r_{n+1}$ of $\gamma$\iftoggle{BW}{}{ (in red)}, with $\tau(r_1) = r_{n+1}$, and parametrize each $r_i$ with $\matR$ via arc-length. We set $\theta' = \theta L(\gamma)/2pi$.}
\label{spezza_su:fig}
\end{center}
\end{figure}

We lift this description to the universal cover $\matH^2$. We fix a lift $l$ of $\bar \eta = \bar{\eta}^0$ and pick $n+1$ consecutive intersections $r_1,\ldots, r_{n+1}$ of $l$ with the lifts of $\gamma$ as in Figure \ref{spezza_su:fig}. The hyperbolic transformation $\tau$ with axis $l$ corresponding to $\eta$ sends $r_1$ to $r_{n+1}$. We parametrise each $r_i$ with $\matR$ via arc length.

The closed geodesic $\bar{\eta}^\theta$, represented as a broken geodesic, lifts to a broken geodesic which starts at some point $x_1-\theta L(\gamma)/2\pi\in r_1$ and arrives at some other point $x_2\in r_2$, then jumps on the left at distance $\theta L(\gamma)/2\pi$ and starts again from $x_2-\theta L(\gamma)/2\pi$, and so on until it reaches the point $\tau(x_1)\in \tau(r_1) = r_{n+1}$. If we make the points $x_1\in r_1, \ldots, x_n \in r_n$ vary we get various broken paths in this way, but only one arrives and exits from each line $r_i$ with the same incidence angles and thus represents $\bar{\eta}^\theta$. The other broken paths represent piecewise-geodesic curves homotopic to $\bar{\eta}^\theta$ and are therefore longer than $\bar{\eta}^\theta$. Hence
$$\ell^\eta(m_\theta^\gamma) = \min \left\{\sum_{i=1}^n d\bigg(x_i-\frac{\theta L(\gamma)}{2\pi}, x_{i+1}\bigg) \ \Bigg|\ (x_1,\ldots, x_n) \in \matR^n \right\}$$
where $x_{n+1} = \tau(x_1)$.
We can now prove that the function $\theta \mapsto \ell^\eta(m_\theta^\gamma)$ is proper and strictly convex. The function
\begin{align*}
\psi\colon \matR^{2n} & \longrightarrow \matR \\
(x_1,y_1, \ldots, x_n,y_n) & \longmapsto \sum_{i=1}^n d(y_i,x_{i+1})
\end{align*}
where $x_{n+1}=\tau(x_1)$ is strictly convex and proper by Proposition \ref{distanza:convessa:prop}. The auxiliary function
\begin{align*}
\phi\colon \matR^{2n}\times \matR & \longrightarrow \matR \\
(x,\theta) & \longmapsto \psi(x)
\end{align*}
is only convex, but its restriction to the subspace
$$H = \bigg\{y_i = x_i - \frac{\theta l(\gamma)}{2\pi} \bigg\}$$
is strictly convex and proper, because the subspace $H$ is not parallel to the direction $(0,\ldots,0,1)$. The coordinates $x_i$ and $\theta$ identify $H$ with $\matR^n \times \matR$. The restriction $f = \phi|_H$ is hence a function $f\colon \matR^n \times \matR \to \matR$ and we obtain
$$\ell^\eta(m_\theta^\gamma) = \min \left\{f(x,\theta) \ \big| \ x \in \matR^n \right\}.$$
By Exercise \ref{convessa:ex} the function $\theta \mapsto \ell^\eta(m_\theta^\gamma)$ is strictly convex and proper. 
\end{proof}

We now employ this convexity property to prove some facts on the earthquakes and the Teichm\"uller space.

\subsection{Earthquakes on essential multicurves}
We use the convexity of the length functions to prove the following. We suppose again that $g\geqslant 2$.

\begin{cor} \label{distinti:cor}
For every simple closed curve $\gamma$, the earthquake action $E_\gamma$ on $\Teich(S_g)$ is free.
\end{cor}
\begin{proof}
Suppose by contradiction that $m = m_{\theta_0}^\gamma$ for some $\theta_0>0$. Then $m= m_{n\theta_0}^\gamma$ for every $n\in\matZ$. Let $\eta$ be a simple closed curve with $i(\eta,\gamma)>0$, which exists by Exercise \ref{eta:prop}; the function $\theta \mapsto \ell^\eta(m_\theta^\gamma)$ is strictly convex and constant on $\{n\theta_0, n \in \matZ\}$, a contradiction.
\end{proof}

The earthquake action is defined more generally for essential multicurves. An essential multicurve $\mu = \gamma_1\sqcup \cdots \sqcup \gamma_k$ of $S_g$ determines an action
\begin{align*}
E_\mu\colon \matR^{k} \times \Teich(S_g) & \longrightarrow \Teich(S_g) \\
 (\theta, m) & \longmapsto m_\theta^{\mu}
\end{align*}
where $\theta = (\theta_1,\ldots, \theta_{k})$ and $m_\theta^\mu = m_{\theta_1}^{\gamma_1} \circ \cdots \circ m_{\theta_{k}}^{\gamma_{k}}$. Note that the actions on disjoint curves commute.

\begin{figure}
\begin{center}
\includegraphics[width = 9cm] {\iftoggle{BW}{dual_curves-BW}{dual_curves}}
\nota{Choose for each component $\gamma_i$ of a pants decomposition $\mu$ a curve $\gamma_i'$ that intersects $\gamma_i$ in one or two points and is disjoint from the other components of $\mu$. There are two cases to consider, depending on whether the two pants adjacent to $\gamma_i$ are distinct (left) or not (right). }
\label{dual_curves:fig}
\end{center}
\end{figure}

\begin{cor}
For every essential multicurve $\mu$, the earthquake action $E_\mu$ on $\Teich(S_g)$ is free. 
\end{cor}
\begin{proof}
We may complete $\mu$ to a pants-decomposition $\mu=\gamma_1\sqcup \ldots \sqcup \gamma_{3g-3}$. Pick for every $i=1,\ldots, 3g-3$ a curve $\gamma_i'$ as in Figure \ref{dual_curves:fig} such that $i(\gamma_i,\gamma_i') > 0$ for all $i$ and $i(\gamma_i, \gamma_j')=0$ for all $i\neq j$. 

Suppose by contradiction that $m=m_\theta^\mu$ for some $\theta\neq 0$: hence $m_{n\theta}^\mu = m$ for all $n\in\matZ$. There is an $i$ such that $\theta_i\neq 0$. The length function $\ell^{\gamma_i'}(m_\theta^\mu)$ depends only on $\theta_i$ and not on the other coordinates of $\theta$: therefore it equals $\ell^{\gamma_i'}(m_{\theta_i}^{\gamma_i})$ which is strictly convex, a contradiction.
\end{proof}

\section{Fenchel--Nielsen coordinates}
It is now time to fix a global set of coordinates for the Teichm\"uller space when $g\geqslant 2$. These are the \emph{Fenchel--Nielsen coordinates} and they identify $\Teich(S_g)$ with $\matR^{6g-6}$, more precisely with $\matR_{>0}^{3g-3} \times \matR^{3g-3}$. 

\subsection{The coordinates}
We want to construct a parametrisation for $\Teich(S_g)$ when $g\geqslant 2$. To identify a finite-dimensional vector space with $\matR^n$ one needs to fix a basis; likewise, here the parametrisation depends on the choice of a  frame.\index{Fenchel--Nielsen coordinates} 

Let $S_g$ be oriented. A \emph{frame} for $S_g$ consists of two essential multicurves $\mu$ and $\nu$ in minimal position, such that:

\begin{enumerate}
\item the multicurve $\mu$ is a pants decomposition,
\item the multicurve $\nu$ decomposes every pair-of-pants in two hexagons.
\end{enumerate}

\begin{figure}
\begin{center}
\includegraphics[width = 10cm] {\iftoggle{BW}{decomposizione2-BW}{decomposizione}}
\nota{A frame for the Fenchel-Nielsen coordinates consists of a (\iftoggle{BW}{light grey}{red}) pants decomposition $\mu$ and a (\iftoggle{BW}{dark grey}{blue}) transverse multicurve $\nu$ that cuts each pair of pants into two hexagons. The number of components of $\mu$ is $3g-3$, that of $\nu$ can vary.}
\label{decomposizione:fig}
\end{center}
\end{figure}

An example that generalises easily to any genus $g\geq 2$ is shown in Figure \ref{decomposizione:fig}. The pants decomposition $\mu = \gamma_1\sqcup \ldots \sqcup \gamma_{3g-3}$ consists of $3g-3$ curves, while the number of curves in $\nu$ is not fixed a priori and depends on our choice of $\nu$.
We now show that a frame induces a \emph{Fenchel--Nielsen} map
\begin{align*}
\FN\colon \Teich(S_g) & \longrightarrow \matR_{>0}^{3g-3} \times \matR^{3g-3} \\
m & \longmapsto  (l_1,\ldots, l_{3g-3}, \theta_1, \ldots, \theta_{3g-3}).
\end{align*}
The map $\FN$ is defined as follows. Let $m\in \Teich(S_g)$ be a hyperbolic metric. The $3g-3$ \emph{length parameters} $l_i = \ell^{\gamma_i}(m)$ are defined using the length functions: the multicurve $\mu$ has a unique geodesic representative 
$$\bar \mu = \bar{\gamma}_1\sqcup \ldots \sqcup \bar{\gamma}_{3g-3}$$ 
in the metric $m$ by Corollary \ref{isotopia:geodetiche:cor}, and $l_i$ is the length of $\bar{\gamma}_i$. Note that these parameters depend only on $\mu$ and not on $\nu$.

The \emph{torsion angles} $\theta_i$ are more subtle to define: the geodesic multicurve $\bar \mu$ decomposes $S_g$ into geodesic pairs-of-pants, and the angle $\theta_i$ measures somehow the way the two geodesic pairs-of-pants are glued along the closed geodesic $\bar \gamma_i$. The precise definition of $\theta_i$ needs the auxiliary multicurve $\nu$. 

\begin{figure}
\begin{center}
\includegraphics[width = 11cm] {\iftoggle{BW}{Teichmuller0-BW}{Teichmuller0}}
\nota{A closed geodesic $\bar \gamma_1$ and the two adjacent pairs-of-pants. The torsion parameter $\theta_1$ measures the distance (in the universal covering) between two orthogeodesics (\iftoggle{BW}{in light grey}{coloured in green}) via the formula $\theta_i = \frac{2\pi s_i}{l_i}$.}
\label{Teichmuller:fig}
\end{center}
\end{figure}

We fix $i=1$ for simplicity and define $\theta_1$.
Figure \ref{Teichmuller:fig}-(left) shows the two geodesic pants adjacent to $\bar{\gamma}_1$ (they might coincide). The second multicurve $\nu$ intersects these pants in four \iftoggle{BW}{dark grey}{blue} arcs, two of which $\lambda$, $\lambda'$ intersect $\bar{\gamma}_1$: we pick one, say $\lambda$. We fix a lift $\tilde P\in\matH^2$ of $P=\bar{\gamma}_1\cap\lambda$ and we lift all the curves incident to $P$: the geodesic $\bar{\gamma}_1$ lifts to a line $\tilde{\gamma}_1$ and $\lambda$ lifts to a (non-geodesic) curve $\tilde\lambda$ that connects two lifts $\tilde{\gamma}_2$ and $\tilde{\gamma}_3$ of the closed geodesics $\bar{\gamma}_2$ and $\bar{\gamma}_3$. See Figure \ref{Teichmuller:fig}-(right).

We draw as in the figure the unique orthogeodesics connecting $\tilde{\gamma}_1$ to $\tilde{\gamma}_2$ and $\tilde{\gamma}_3$ and we denote by $s_1$ the signed length of the segment in $\tilde{\gamma}_1$ comprised between these two orthogeodesics, with positive sign if (as in the figure) an observer walking on a orthogeodesic towards $\tilde{\gamma}_1$ sees the other orthogeodesic on its left (here we use the orientation of $S_g$). 

By repeating this construction for each $\bar{\gamma}_i$ we find some real numbers $s_i$. Finally, the torsion parameter $\theta_i$ is 
$$\theta_i = \frac {2\pi s_i}{l_i}.$$

\begin{figure}
\begin{center}
\includegraphics[width = 5cm] {\iftoggle{BW}{Teichmuller2-BW}{Teichmuller2}}
\nota{If we pick $\lambda'$ instead of $\lambda$ we find a segment of the same length $s_1$. This holds because the two geodesic pairs-of-pants incident to $\gamma_1$ decompose into two hexagons isometric to $A$ and $B$ as shown, and the sides of $A$ and $B$ contained in $\tilde{\gamma}_1$ have the same length $\frac{l(\gamma_1)}2$, that is half the length of $\gamma_1$.}
\label{Teichmuller2:fig}
\end{center}
\end{figure}

\begin{teo}[Fenchel-Nielsen coordinates]
The map $\FN$ is well-defined and is a bijection.
\end{teo}
\begin{proof}
We first note that in the definition of the torsion parameters we could have chosen $\lambda'$ instead of $\lambda$. We would have obtained the same length $s_i$ as shown in Figure \ref{Teichmuller2:fig}. Moreover a hyperbolic metric $m'$ isometric to $m$ through a diffeomorphism $\varphi$ isotopic to the identity has the same parameters $l_i$ and $\theta_j$ since they depend only on the isotopy classes of $\mu$ and $\nu$. Therefore $\FN$ is well-defined.

We prove that $\FN$ is surjective. For every vector $(l_1,\ldots,l_{3g-3})\in\matR_{>0}^{3g-3}$ we may use Proposition \ref{pantaloni:prop} and construct a metric on $S_g$ by assigning to each pair-of-pants of the pants decomposition $\mu$ the (unique) hyperbolic metric with boundary lengths $l_i$. We get a metric with some arbitrary torsion angles $\theta$, which can be changed arbitrarily by an earthquake along $\mu$: it is easy to check that an earthquake with angles $\theta'$ changes the torsion angles from $\theta$ to $\theta+\theta'$, hence any torsion parameter can be realised and $\FN$ is surjective.

We prove that $\FN$ is injective. If $\FN(m) = \FN(m')$, up to acting via earthquakes we may suppose that $\FN(m)=\FN(m') = (l_1,\ldots,l_{3g-3},0,\ldots,0)$. Since the torsion parameter is zero, the orthogeodesics in Figure \ref{Teichmuller:fig}-(right) match and project in $S_g$ to a geodesic multicurve $\bar \nu$ isotopic to $\nu$ and orthogonal to $\bar\mu$. Therefore $S_g\setminus (\bar\mu \cup \bar\nu)$ is a tessellation of $S_g$ into right-angled hexagons, determined by the lengths $l_i$. Both metrics $m$ and $m'$ have the same tessellation and are hence isometric, via an isometry which is isotopic to the identity.
\end{proof}

\begin{oss}
As shown in the proof, the torsion parameters for $m$ are zero if and only if the geodesic representatives $\bar \nu$ and $\bar \mu$ of $\nu$ and $\mu$ are everywhere orthogonal. 
\end{oss}


\subsection{Length functions of $9g-9$ curves}
It is now natural to ask whether the length functions determine every point in $\Teich(S_g)$. The answer is positive; as usual, to warm up we start by examining the torus.

\begin{figure}
\begin{center}
\includegraphics[width = 7cm] {\iftoggle{BW}{torus_3_curves-BW}{torus_3_curves}}
\nota{The curves $\gamma$\iftoggle{BW}{}{ (red)}, $\gamma'$\iftoggle{BW}{}{ (blue)}, and $\gamma'' = T_{\gamma}(\gamma')$\iftoggle{BW}{}{ (green)} on the torus.}
\label{torus_3_curves:fig}
\end{center}
\end{figure}

Let $\gamma,\gamma'$ be two simple closed curves in the torus $T$ with $i(\gamma, \gamma')=1$ and let $\gamma'' = T_{\gamma}(\gamma')$ be obtained by Dehn twisting $\gamma'$ along $\gamma$, see Figure \ref{torus_3_curves:fig}.

\begin{prop} \label{3:prop}
The map
\begin{align*}
L\colon \Teich(T) & \longrightarrow \matR_{>0}^{3} \\
m & \longmapsto \big(\ell^{\gamma}(m), \ell^{\gamma'}(m), \ell^{\gamma''}(m)\big)
\end{align*}
is injective.
\end{prop}
\begin{proof}
After fixing $\gamma$ and $\gamma'$ as a homology basis, we have $\gamma = (1,0)$, $\gamma' = (0,1)$, and $\gamma'' = (1,-1)$. Proposition \ref{lunghezza:toro:prop} gives
$$L(z) = \left(\frac{1}{\sqrt{\Im z}}, \frac{|z|}{\sqrt{\Im z}}, \frac{|z-1|}{\sqrt{\Im z}} \right)$$
which is easily seen to be injective on $\Teich(T) = H^2$.
\end{proof}

A similar set of $9g-9$ curves does the job on $S_g$ when $g\geqslant 2$.
Let $\mu = \gamma_1\sqcup \ldots \sqcup \gamma_{3g-3}$ be a pants decomposition for $S_g$. For each $\gamma_i$ we choose a curve $\gamma_i'$ as in Figure \ref{dual_curves:fig}, and we indicate by $\gamma_i'' = T_{\gamma_i}(\gamma_i')$ the curve obtained by Dehn-twisting $\gamma_i'$ along $\gamma_i$, see an example in Figure \ref{dual_curves2:fig}.

\begin{figure}
\begin{center}
\includegraphics[width = 4cm] {\iftoggle{BW}{dual_curves2-BW}{dual_curves2}}
\nota{The curves $\gamma_i$\iftoggle{BW}{}{ (red)}, $\gamma_i'$\iftoggle{BW}{}{ (blue)}, and $\gamma_i'' = T_{\gamma_i}(\gamma_i')$\iftoggle{BW}{}{ (green)} when $\gamma_i$ is adjacent twice to the same pair of pants. It is an instructive exercise to draw $\gamma_i'$ and $\gamma_i''$ when $\gamma_i$ is incident to distinct pair of pants.}
\label{dual_curves2:fig}
\end{center}
\end{figure}

\begin{prop} \label{9g-9:prop}
The map
\begin{align*}
L\colon \Teich(S_g) & \longrightarrow \matR_{>0}^{9g-9} \\
m & \longmapsto \big(\ell^{\gamma_i}(m), \ell^{\gamma_i'}(m), \ell^{\gamma_i''}(m)\big)
\end{align*}
is injective.
\end{prop}
\begin{proof}
We compose $L$ with $\FN^{-1}$ and obtain a map
\begin{align*}
L\circ \FN^{-1}\colon \matR^{3g-3}_{>0} \times \matR^{3g-3} & \longrightarrow \matR_{>0}^{9g-9} \\
(l_i, \theta_i) \quad \ \ & \longmapsto (l_i, l_i', l_i'') 
\end{align*}
We prove that it is injective: it suffices to consider the case where the values $l_i$ are fixed and $\theta_i$ vary. Note that $\gamma_i'$ and $\gamma_i''$ intersect $\gamma_j$ if and only if $i=j$: hence $l_i'$ and $l_i''$ depend only on $\theta_i$ and not on the other torsion parameters $\theta_j$. Proposition \ref{strettamente:convessa:prop} says that $l_i' = f(\theta_i)$ is strictly convex and Proposition \ref{root:prop} gives $l_i'' = f(\theta_i+2\pi)$. 
A strictly convex proper function $f\colon\matR \to \matR$ is at most 2 to 1, hence the function 
\begin{align*}
\matR & \longrightarrow \matR \times \matR \\
\theta_i & \longmapsto \big(f(\theta_i), f(\theta_i + 2 \pi)\big)
\end{align*}
is injective. Therefore $L$ is injective. 
\end{proof}

\subsection{Collar lemma}
The thick-thin decomposition theorem implies that the closed geodesics of length smaller than a Margulis constant $\varepsilon_2$ on a complete hyperbolic surface are simple and have disjoint $R$-neighbourhoods (see Corollary \ref{simple:small:geodesics:cor}). When the curves are very short, one may choose $R$ to be very large: this fact is called the \emph{collar lemma} and we prove it directly using elementary tools.\index{collar lemma}

\begin{figure}
\begin{center}
\includegraphics[width = 11cm] {\iftoggle{BW}{collar-BW}{collar}}
\nota{Pick a segment of length $l$ and draw two perpendiculars at the endpoints: this determines a quadrilateral with two ideal vertices; let $f(l)$ be the distance between the opposite sides $l$ and $r$ (left). A geodesic pair-of-pants is the union of two isometric hexagons, which form in $\matH^2$ a right-angled octagon as drawn. The picture shows that the closed boundary geodesics $a$ and $b$ have disjoint $f(a)$ and $f(b)$-neighbourhoods, coloured here in \iftoggle{BW}{light grey}{yellow} (right).} 
\label{collar:fig}
\end{center}
\end{figure}

For any number $l>0$, draw the quadrilateral as in Figure \ref{collar:fig}-(left) and define $f(l)$ to be the distance between its opposite sides $l$ and $r$.

\begin{ex} The function $f\colon \matR_{>0}\to \matR_{>0}$ is strictly decreasing and a homeomorphism. In particular we get
$\lim_{l \to 0} f(l) = \infty$. 
Explicitly, we have 
$$\sinh f(l) = \frac{1}{\sinh \frac l2}.$$ 
\end{ex}
\begin{proof}[Hint]
Put $l$ in vertical position in $H^2$ and use Lemma \ref{cosh:cos:lemma}.
\end{proof}

The function $f$ is simple to define, and is particularly useful.

\begin{prop} \label{collar:pants:prop}
Let $P$ be a geodesic pair-of-pants with boundary lengths $a,b$, and $c$. The $f(a)$, $f(b)$, and $f(c)$-neighbourhoods of the boundary components form three disjoint collars.
\end{prop}
\begin{proof}
Consider two boundary components $a$ and $b$. 
The geodesic pair-of-pants $P$ is divided into two isometric hexagons, and we lift them to $\matH^2$ where they form a right-angled octagon as in Figure \ref{collar:fig}-(right). The picture shows that the $f(a)$ and $f(b)$-neighbourhoods of $a$ and $b$ are disjoint.
\end{proof}

Here is the collar lemma.

\begin{lemma}[Collar lemma] \label{collar:lemma}
Let $g\geqslant 2$ and $S_g$ have a hyperbolic metric. Disjoint simple closed geodesics $\gamma_1, \ldots, \gamma_k$ of length $l_1,\ldots, l_k$ have disjoint tubular $f(l_i)$-neighbourhoods. 
\end{lemma}
\begin{proof}
We may suppose (by adding more simple closed geodesics if necessary) that the closed geodesics form a pants decomposition, and it suffices to consider two curves that cobound the same pair-of-pants. Proposition \ref{collar:pants:prop} applies.
\end{proof}

We recall from Corollary \ref{simple:small:geodesics:cor} that every closed geodesic shorter than $\varepsilon_2$ is simple. Therefore very short closed geodesics are simple and have large disjoint tubular neighbourhoods: the more we shrink the curves, the larger are their neighbourhoods, and hence the larger is the diameter of the surface (the diameter of a metric space is the supremum of the distance of its points). In particular, if on a sequence of closed hyperbolic surfaces the injectivity radius tends to zero, their diameters must tend to infinity.

Among the many consequences of the collar lemma, we focus on a simple inequality which relates the geometric intersection of simple closed geodesics to their lengths. We denote by $L(\gamma)$ the length of $\gamma$.

\begin{cor}  \label{strizza:cor}
Let $g\geqslant 2$ and $S_g$ have a hyperbolic metric. Let $\gamma$ and $\eta$ be two simple closed geodesics in $S_g$. The following inequality holds:
$$L(\eta)\geqslant 2i(\eta, \gamma) \cdot f(L(\gamma)).$$
\end{cor}
\begin{proof}
The geodesic $\gamma$ has a tubular $f(L(\gamma))$-neighbourhood. The geodesic $\eta$ intersects $\gamma$ in $i(\eta,\gamma)$ points and hence crosses the tubular neighbourhood at least $i(\eta,\gamma)$ times, each with a segment of length $\geqslant 2f(L(\gamma))$.
\end{proof}

\subsection{A topology for the Teichm\"uller space} \label{Topology:Teich:subsection}
There are various equivalent ways to assign a topology to the Teichm\"uller space. 
On the torus $T$, we have seen that $\Teich(T)$ can be identified with the hyperbolic plane $\matH^2$, and the mapping class group acts as isometries on it: we could not hope for a better picture of $\Teich(T)$ and we are fully satisfied.

When $g\geqslant 2$ we could similarly use the Fenchel-Nielsen coordinates and give $\Teich(S_g)$ the topology of $\matR^{6g-6}$, but then to be honest we should also check that the topology does not depend on the frame... we prefer to equip the Teichm\"uller space with an intrinsic topology and then prove that the Fenchel-Nielsen coordinates are homeomorphisms.\index{$\calS$} 

We indicate by $\calS= \calS(S_g)$ the set of all the non-trivial simple closed curves in $S_g$, considered up to isotopy and orientation reversal (we say that the curves are \emph{unoriented}). Each element $\gamma\in\calS$ induces a length function
$$\ell^\gamma\colon \Teich(S_g) \longrightarrow \matR_{>0}.$$
We indicate as usual with $\matR^\calS$ the set of all functions $\calS \to \matR$ and give it the usual product topology (the weakest one such that all the projections are continuous). The natural map
\begin{align*}
\Teich(S_g) & \longrightarrow \matR^\calS \\
m & \longmapsto \big( \gamma \longmapsto \ell^\gamma(m) \big)
\end{align*}
is injective by Propositions \ref{3:prop} and \ref{9g-9:prop}. We may hence consider $\Teich(S_g)$ as a subspace of $\matR^\calS$ and assign it the subspace topology. This topology on $\Teich(S_g)$ is the weakest one where the length functions $\ell^\gamma$ are continuous.
Recall that a topological space is \emph{second-countable} if it has a countable base.

\begin{prop}
The space $\matR^\calS$ is Hausdorff and second-countable.
\end{prop}
\begin{proof}
Every product of Hausdorff spaces is Hausdorff, and every countable product of second-countable spaces is second-countable.
\end{proof}

We recall the following topological fact.

\begin{prop} \label{propria:chiusa:prop}
Let $f\colon X \to Y$ be a continuous and proper map between topological spaces. If $Y$ is Hausdorff and second-countable then $f$ is closed.
\end{prop}

Proper maps onto reasonable spaces are closed. If they are also injective, we can obtain more.

\begin{cor} \label{propria:chiusa:cor}
Let $f\colon X \to Y$ be a continuous, proper, and injective map between topological spaces. If $Y$ is Hausdorff and second-countable then $f$ is a homeomorphism onto its image. 
\end{cor}

We will use this corollary in a moment.
Recall that every isometry $\varphi\in \Iso^+(H^2) = \PSLR$ has a trace $\tr \varphi$ defined only up to sign, whose modulus is $>2$ precisely when $\varphi$ is hyperbolic, see Proposition \ref{PSLR:prop}.

\begin{prop} \label{trace:prop}
Let $S = \matH^2/_\Gamma$ be an orientable hyperbolic surface. Every hyperbolic transformation $\varphi\in\Gamma$ produces a closed geodesic $\gamma$ in $S$
with
$$|\tr\varphi| = 2\cosh \frac {L(\gamma)} 2.$$
\end{prop}
\begin{proof}
Up to conjugacy we have $\varphi(z) = e^{L(\gamma)} z$. The matrix is 
$$\varphi = \left(\begin{matrix} e^{\frac{L(\gamma)}2} & 0 \\ 0 & e^{-\frac{L(\gamma)}2} \end{matrix}\right)$$
hence $|\tr\varphi| = 2\cosh \frac{L(\gamma)}2$.
\end{proof}

In particular, the length of $\gamma$ depends continuously on the transformation $\varphi$. We will use this fact to prove the following. We suppose $g\geqslant 2$.

\begin{prop} \label{FN:omeom:prop}
The Fenchel-Nielsen map 
$$\FN\colon \Teich(S_g) \longrightarrow \matR_{>0}^{3g-3}\times\matR^{3g-3}$$
is a homeomorphism.
\end{prop}
\begin{proof}
We consider $\Teich(S_g)$ inside $\matR^\calS$ and examine the inverse map
$$\FN^{-1}\colon\matR_{>0}^{3g-3}\times\matR^{3g-3} \longrightarrow \matR^\calS.$$

We prove that $\FN^{-1}$ is continuous. The map $\FN^{-1}$ assigns to the parameters $(l_i, \theta_i)$ a metric on $S_g$ constructed by attaching right-angled hexagons. Both the hexagons and the attaching maps depend continuously on the parameters $(l_i, \theta_i)$ and lift to a tessellation of $\matH^2$ into hexagons. Since the decomposition into hexagons varies continuously, its deck transformations vary continuously in $\PSLR$ and hence the length functions too by Proposition \ref{trace:prop}. Therefore $\FN^{-1}$ is continuous.

We prove that $\FN^{-1}$ is proper. Take a diverging sequence of parameters $(l_i,\theta_i)$ (that is, without converging subsequences) in $\matR_{>0}^{3g-3}\times \matR^{3g-3}$: we need to show that its image is also a diverging subsequence. The thesis is equivalent to show that the length function of some curve goes to infinity. If $l_i\to +\infty$ for some $i$ we are done. If $l_i\to 0$, the length of any curve intersecting essentially the shrinking curve $\gamma_i$ goes to infinity by Corollary \ref{strizza:cor}. It remains to consider the case where the length parameters $l_i$ converge to some non-zero value, but some twist parameter $\theta_j$ goes to infinity: in that case the length of any curve intersecting the twisted curve $\gamma_j$ goes to infinity by Proposition \ref{strettamente:convessa:prop}.

Finally, the map $\FN^{-1}$ is a homeomorphism onto its image by Corollary \ref{propria:chiusa:cor}. The proof is complete.
\end{proof}

During the proof we have also discovered the following.

\begin{prop} \label{diverges:infinity:prop}
If a sequence $m_i\in\Teich(S_g)$ diverges, there is a $\gamma \in \calS$ such that $\ell^\gamma(m_i) \to \infty$ on a subsequence.
\end{prop}

Recall that the action of a topological group $G$ on a topological space $X$ is \emph{continuous} if the action map $G\times X \to X$ is continuous. This implies that $G$ acts on $X$ by homeomorphisms. We give the mapping class group $\MCG(S_g)$ the discrete topology.

\begin{prop}
The earthquakes and mapping class group actions on the Teichm\"uller space are continuous.
\end{prop}
\begin{proof}
The mapping class group acts on $\calS$ by permutations, hence its action on $\matR^\calS$ is continuous. On Fenchel-Nielsen coordinates the earthquake action sends $\theta$ to $\theta+\theta'$ and is hence continuous.
\end{proof}

The immersion in $\matR^{9g-9}$ is also a topological embedding.

\begin{prop}
The injective representation $\Teich(S_g) \hookrightarrow \matR^{9g-9}$
furnished by Proposition \ref{9g-9:prop} is a homeomorphism onto its image.
\end{prop}
\begin{proof}
Using Fenchel-Nielsen coordinates the map is clearly continuous. The proof that it is proper is as in Proposition \ref{FN:omeom:prop}.
\end{proof}

\subsection{Surfaces of finite type} \label{Teich:finite:type:subsection}
We have considered only closed surfaces $S_g$ for simplicity, but most of the arguments exposed in this chapter extend easily to all surfaces $S_{g,b,p}$ of finite type with negative Euler characteristic. 

The Teichm\"uller space $\Teich(S_{g,b,p})$ is the set of all the complete hyperbolic metrics with geodesic boundary, considered up to isometries that are isotopic to the identity. Fenchel-Nielsen coordinates are defined analogously: the surface decomposes into $-\chi(S_{g,b,p}) = 2g+b+p-2$ pairs-of-pants, and the interior curves of the decomposition are 
$$\frac 12 \big(3(2g+b+p-2) - b - p \big) = 3g+b+p-3.$$
The Fenchel-Nielsen coordinates are
$$\big(l_1, \ldots, l_{3g+b+p-3}, l_1^\partial, \ldots, l_b^\partial, \theta_1,\ldots, \theta_{3g+b+p-3}\big)$$
where the $l_i$ and $l_j^\partial$ are the length parameters of the $3g+b+p-3$ interior and $b$ boundary curves, and the $\theta_i$ are the torsion angles of the interior curves. We get a bijection
$$\FN\colon \Teich (S_{g,b,p}) \longrightarrow \matR_{>0}^{3g+2b+p-3} \times \matR^{3g+b+p-3}$$
which is a homeomorphism with respect to the natural topology on Teichm\"uller space as a subset of $\matR^\calS$. Therefore the Teichm\"uller space is homeomorphic to a ball of dimension $-3\chi(S_{g,b,p})-p$.

For instance, the Teichm\"uller space of a pair-of-pants is $\matR_{>0}^3$ parametrized by the lengths of the boundary geodesics, while that of a thrice-punctured sphere is a point (there is a unique metric).

An alternative description of the Teichm\"uller space for punctured surfaces, with ideal triangles playing the role of pairs-of-pants, is described in the next section.

\section{Shear coordinates}
The Teichm\"uller space of a punctured surface may also be parametrized using ideal triangles instead of pairs-of-pants: this viewpoint is maybe a bit simpler, and generalises successfully to dimension three (via ideal tetrahedra). The coordinates that it produces are called \emph{shear coordinates}.\index{shear coordinates}

\subsection{Ideal triangulations}
In dimension two and three it is customary to relax the definition of \emph{triangulation}, originally restricted to simplicial complexes, see Section \ref{triangulations:subsection}. We prefer to define a triangulation in a looser sense, as a finite set of triangles glued together by pairing their edges. 

Let $\Delta_1,\ldots, \Delta_{2k}$ be an even number of identical copies of the standard oriented 2-simplex. A \emph{triangulation} $\calT$ is a partition of the $6k$ edges of the triangles into $3k$ pairs, and for each pair a simplicial isometry between the two edges. The triangulation is \emph{oriented} if the simplicial isometries are orientation-reversing. If we glue the triangles along the isometries we get a compact surface $S$: we always suppose that $S$ is connected and $\calT$ is oriented, hence $S = S_g$ for some $g\geqslant 0$.

The surface $S$ is triangulated with $2k$ triangles, $3k$ edges, and some $p$ vertices. Vertices, edges, and triangles form a cellularisation of $S$, but not necessarily a simplicial complex: for instance two or three vertices of the same triangle $\Delta_i$ can be identified to a single one along the process, as the following exercise shows.

\begin{ex}
Construct a triangulation of the torus with one vertex, three edges, and two triangles.
\end{ex}

Let $\Sigma$ be the non-compact surface obtained by removing the $p$ vertices of the triangulation $\calT$ from $S$: we say that $\calT$ is an \emph{ideal triangulation} for $\Sigma$. The surface $\Sigma$ is a \emph{punctured surface}, \emph{i.e.}~$\Sigma = S_{g,0,p}$ with $p\geqslant 1$.\index{ideal triangulation}

\begin{figure}
\begin{center}
\includegraphics[width = 9 cm] {\iftoggle{BW}{triangulated_surface-BW}{triangulated_surface}}
\nota{The standard representation of $S_g$ by identifying the edges of a $4g$-gon has only one vertex; if we triangulate the $4g$-gon we get an ideal triangulation for $S_{g,0,1}$ (left). A move that modifies a triangulation by adding a new vertex (right).}
\label{triangulated_surface:fig}
\end{center}
\end{figure}

\begin{prop} We have $\chi(\Sigma) = -k <0$. Every punctured surface of negative Euler characteristic has an ideal triangulation.
\end{prop}
\begin{proof}
We have $\chi(\Sigma) = \chi(S) -p = p-3k+2k-p = -k$. An ideal triangulation for $S_{0,0,3}$ is constructed by attaching $\Delta_1$ and $\Delta_2$ via the obvious maps. An ideal triangulation for $S_{g,0,1}$ is constructed in Figure \ref{triangulated_surface:fig}-(left) for all $g\geqslant 1$. An ideal triangulation for $S_{g,0,p}$ with $\chi(S_{g,0,p})<0$ and $p\geqslant 2$ is obtained from one of these by increasing $p$ as in Figure \ref{triangulated_surface:fig}-(right).
\end{proof}

\subsection{Hyperbolic ideal triangulations}
Let $\calT$ be an oriented triangulation with triangles
$\Delta_1,\ldots, \Delta_{2k}$. We substitute every $\Delta_i$ with an ideal hyperbolic triangle (recall that it is unique up to isometry) and pair the edges with some orientation-reversing isometries. The resulting oriented punctured surface $\Sigma$ now inherits a hyperbolic structure of area $k\pi$, since an ideal triangle has area $\pi$. We call $\calT$ a \emph{hyperbolic ideal triangulation} for $\Sigma$.


\begin{figure}
\begin{center}
\includegraphics[width = 9 cm] {\iftoggle{BW}{triangoli_ideali-BW}{triangoli_ideali}}
\nota{On an ideal triangle every edge has a unique perpendicular which ends at the opposite vertex. The base of this perpendicular is the \emph{midpoint} of the edge and the three perpendiculars intersect in a point called \emph{barycenter} (left). There is a unique configuration of three pairwise tangent horocycles centred at the ideal vertices, and their tangency points are the midpoints of the edges (right).}
\label{triangoli_ideali:fig}
\end{center}
\end{figure}

\begin{figure}
\begin{center}
\includegraphics[width = 11 cm] {\iftoggle{BW}{shear-BW}{shear}}
\nota{The \emph{shear} is the signed distance $d$ between two midpoints after the gluing. By convention the sign is positive $d>0$ if the points are as shown here, that is an observer that arrives at the midpoint of an edge sees the other midpoint on its left: we need the orientation of $\calT$ here (left). We develop a triangulation around a puncture $v$: here $h=3$ (right).}
\label{shear:fig}
\end{center}
\end{figure}

The edges are infinite lines and the gluing isometry is not unique: indeed there is a 1-parameter family of isometries to choose from at every pair of edges, and the hyperbolic structure depends on that -- although the topology of $\Sigma$ does not. To encode this dependence, we note that every edge of an ideal triangle has a preferred \emph{midpoint} defined in Figure \ref{triangoli_ideali:fig}. The orientation-reversing isometry between two edges of two triangles is determined by the signed distance $d$ (called \emph{shear}) of their two midpoints as shown in Figure \ref{shear:fig}-(left). The hyperbolic structure on $\Sigma$ is determined by the \emph{shear coordinates} $d=(d_1, \ldots, d_{3k})$ of the $3k$ edges of the triangulation.

\subsection{Complete solutions}
The hyperbolic structure on $\Sigma$ may be incomplete! Recall that a hyperbolic surface is complete if and only if its universal cover is $\matH^2$, so we now look at the universal cover of $\Sigma$. 

\begin{figure}
\begin{center}
\includegraphics[width = 4cm] {\iftoggle{BW}{universale-BW}{universale}}
\nota{The universal covering of $\Sigma$ is isometric to the interior of a convex set delimited by some lines in $\matH^2$, which is the intersection of (possibly infinitely many) half-planes bounded by disjoint lines.}
\label{universal:fig}
\end{center}
\end{figure}

\begin{prop}
The universal cover $\tilde \Sigma$ of $\Sigma$ is isometric to
the interior of the intersection of some half-planes in $\matH^2$ with disjoint boundary lines (see Figure \ref{universal:fig}). The surface $\Sigma$ is complete if and only if $\tilde \Sigma \isom\matH^2$.
\end{prop}

\begin{proof}
Recall from Section \ref{non-complete:subsection} that there is a developing map $D\colon \tilde {\Sigma} \to \matH^2$. We prove that $D$ is injective. 

We can construct $D$ as follows: the hyperbolic ideal triangulation $\calT$ of $\Sigma$ lifts to a hyperbolic ideal triangulation $\tilde \calT$ of $\tilde{\Sigma}$ with infinitely many triangles. Send a triangle of $\tilde \calT$ to an arbitrary ideal triangle of $\matH^2$ and then develop the map $D$ by attaching subsequently all the triangles of $\tilde \calT$.  At each step the image of $D$ is an ideal polygon and we attach a new triangle to the exterior of some side of the polygon. Therefore at each step $D$ is injective and hence $D$ is globally injective.

Since $D$ is injective we identify $\tilde{\Sigma}$ with its open image in $\matH^2$. Pick a point $x\in\partial\tilde{\Sigma}$. Every neighbourhood of $x$ must intersect infinitely many triangles of $\tilde \calT$, hence there is a sequence of edges $e_i$ of $\tilde \calT$ such that $d(e_i,x)$ is monotone decreasing and tends to zero. Since the $e_i$ are disjoint, their endpoints in $\partial \matH^2$ are monotone and tend to two points, which determine a line $l$. We must have $d(l,x)=0$ and hence $x\in l$ and $l\subset \partial \tilde \Sigma$. 

We have discovered that $\partial \tilde \Sigma$ consists of disjoint lines. This proves that the closure of $\tilde\Sigma$ is the intersection of half-planes. A surface $\Sigma$ is complete if and only if $\tilde \Sigma$ is complete, and hence isometric to $\matH^2$.
\end{proof}

The hyperbolic structure on $\Sigma$ depends on the shearing coordinates $d\in \matR^{3k}$, and we now want to understand which values of $d$ produce a complete metric. An ideal vertex $v$ of the triangulation $\calT$ is adjacent to some $h$ triangles, which we denote by $\Delta_{1},\ldots, \Delta_{h}$ for simplicity although there might be repetitions, and to $h$ edges with some shearing coordinates that we also indicate by $d_1,\ldots, d_h$ for simplicity. The vertex $v$ is a puncture of $\Sigma$ and we let $N(v)$ be a small topological punctured closed disc around $v$.

\begin{prop}
The punctured disc $N(v)$ is complete if and only if $d_1 + \ldots + d_h = 0$.
\end{prop}
\begin{proof}
We construct a portion of the developing map $D\colon \tilde \Sigma \to \matH^2$ as follows: we use the half-plane model and send $\Delta_1$ to $\matH^2$ with $v=\infty$ as in Figure \ref{shear:fig}-(right), then we develop the triangulation horizontally. In the picture we have $h=3$ so we reproduce $\Delta_1, \Delta_2, \Delta_3$ and then another copy $\Delta_1'$ of $\Delta_1$. Let $\varphi\in\PSLR$ be the holonomy isometry that sends $\Delta_1$ to $\Delta_1'$.

If $d_1+\ldots + d_h=0$ the midpoints of $\Delta_1$ and $\Delta_1'$ are at the same height as in Figure \ref{shear:fig}-(right) and therefore $\varphi$ is a parabolic map $\varphi\colon z \mapsto z+b$ for some $b\in \matR$. Therefore $N(v)$ is a truncation of the cusp $\matH^2/_{\langle \varphi\rangle }$, which is complete.

\begin{figure}
\begin{center}
\includegraphics[width = 11 cm] {\iftoggle{BW}{shear2-BW}{shear2}}
\nota{When $d_1+\ldots + d_h \neq 0$, the isometry $\varphi$ sending $\Delta_1$ to $\Delta_1'$ is hyperbolic, with some axis $l$ which (up to translations) has endpoints $0$ and $\infty$. Hence $\varphi(z)=e^d z$ and the triangulation develops towards $l$ as shown here.}
\label{shear2:fig}
\end{center}
\end{figure}

If $d_1+\ldots + d_h \neq 0$ the midpoints of $\Delta_1$ and $\Delta_1'$ are at different heights and therefore $\varphi$ is not parabolic. Since $\infty \in \Fix(\varphi)$, the isometry $\varphi$ is a hyperbolic transformation having some other fixed point in $\matR$. Up to translating everything we may suppose that $\Fix(\varphi)=\{0,\infty\}$, hence the axis $l$ of $\varphi$ is the vertical coordinate axis, and $\varphi(z) =  e^d z$ with $d = d_1+\ldots + d_h$. The map $D$ develops towards $l$ as shown in Figure \ref{shear2:fig}: we get $l\subset \partial \tilde\Sigma$, so $\tilde\Sigma \neq\matH^2$ is not complete. 
\end{proof}
 
Summing up, we have a \emph{completeness equation} of type $d_{i_1}+\ldots + d_{i_h}=0$ for each of the $p$ ideal vertices of the triangulation.

\begin{cor}
The hyperbolic structure on $\Sigma$ is complete if and only if $d$ satisfies the $p$ completeness equations.
\end{cor}
\begin{proof}
Let $N(v)$ be a small punctured closed disc at $v$ for every ideal vertex $v$. The closure of $\Sigma \setminus \cup_v N(v)$ is compact. Therefore $\Sigma$ is complete if and only if each $N(v)$ is.
\end{proof} 

\subsection{Shear coordinates for Teichm\"uller space}
The solution space of the completeness equations is some linear subspace $V \subset \matR^{3k}$ of dimension at least $3k-p = -3\chi(\Sigma)-p$. Note that this is precisely the dimension of $\Teich(\Sigma)$, see Section \ref{Teich:finite:type:subsection}. We have constructed a map
$${\rm sh}\colon V \longrightarrow \Teich(\Sigma).$$
We will prove below that $\sh$ is a homeomorphism and hence $\dim V = 3k-p$. We will use the following.

\begin{prop} \label{straightens:prop}
Let $\Sigma$ be a complete hyperbolic punctured surface.
Every ideal triangulation is isotopic to a unique hyperbolic ideal triangulation. 
\end{prop}
\begin{proof}
We describe a procedure that \emph{straightens} every ideal triangulation $\calT$ of $\Sigma = \matH^2/_\Gamma$ to a hyperbolic one. 

Truncate the cusps of $\Sigma$ to get a compact sub-surface. The ideal triangulation $\calT$ is now truncated to a decomposition into hexagons, each with three boundary and three interior edges. We lift the truncated cusps and the hexagons to $\matH^2$. Truncated cusps lift to horoballs, bounded by horocycles, and hexagons lift to hexagons (the hexagons are not geodesic, only smooth). 

An interior edge of a hexagon in $\matH^2$ cannot have both its endpoints in the same horocycle, since it would form a bigon with it, and a bigon does not decompose into hexagons (by a Euler characteristic argument). Therefore every interior edge connects two distinct horocycles, centred at two distinct points of $\partial \matH^2$. Replace this edge with the geodesic line connecting these two ideal points.

If we do this at every interior edge we get a $\Gamma$-invariant hyperbolic ideal triangulation of $\matH^2$, which projects to a hyperbolic ideal triangulation for $\Sigma$ isotopic to $\calT$.  
\end{proof}

We can now parametrize $\Teich(\Sigma)$.
 
\begin{cor} The map $\sh$ is a bijection.
\end{cor}
\begin{proof}
Proposition \ref{straightens:prop} shows that in every metric the initial triangulation $\calT$ straightens to an ideal hyperbolic one and is hence realised by some $d$, so $\sh$ is surjective. Moreover the straightened triangulation is unique and the parameters $d_i$ are obtained intrinsically from it, so $\sh$ is also injective.
\end{proof}

\subsection{Incomplete metrics} \label{incomplete:metrics:subsection}
The shear coordinates $d$ may define an incomplete metric on $\Sigma$, and we now determine its metric completion $\overline \Sigma$.

Recall that every ideal vertex $v$ has a \emph{total shear} $d_v = d_{i_1}+\ldots +d_{i_h}$ which is zero precisely when $\Sigma$ is complete near $v$. 

\begin{figure}
\begin{center}
\includegraphics[width = 11 cm] {\iftoggle{BW}{shear3-BW}{shear3}}
\nota{Let $s$ be a Euclidean line at hyperbolic distance $R$ from $l$. If $R$ is sufficiently small, the line $s$ does not intersect the lower edges of $\Delta_1, \Delta_2, \Delta_3$. Let $S$ be the region lying above $s$.}
\label{shear3:fig}
\end{center}
\end{figure}


\begin{prop}  \label{completion:surface:prop}
The metric completion $\overline\Sigma$ is a hyperbolic surface with geodesic boundary, whose boundary lengths are the absolute values of the total shears of the vertices. 
\end{prop}
\begin{proof}
Let $N(v)$ be a small closed punctured disc around $v$.
If $d_v=0$ we already know that $N(v)$ is a truncated cusp. If $d_v\neq 0$ we develop $N(v)$ in $\matH^2$ as in Figure \ref{shear3:fig}. The hyperbolic transformation is $\varphi(z) = e^{d_v} z$ and the points at distance $R$ from $l$ form a Euclidean line $s$ starting from $0$. Pick $R$ small so that $s$ does not intersect the lower edges of $\Delta_1,\ldots,\Delta_h$. Since $s$ is $\varphi$-invariant, it is contained in the union of all the developed triangles of Figure \ref{shear3:fig} and hence in $\tilde\Sigma$.

The line $s$ projects to a (non-geodesic) circle in $\Sigma$ around $v$, and we can suppose that $\partial N(v)$ is that circle. We have $N(v) = S/_\gamma$ where $S\subset H^2$ is the region lying above $s$. The completion $\overline S $ equals $S\cup l$ and therefore $\overline{N(v)} = \overline S/_\gamma$ is an annulus with two boundary components: the original non-geodesic $\partial N(v)$ and a new geodesic $l/_\gamma$, a curve of length $|d_v|$.
\end{proof}

\begin{figure}
\begin{center}
\includegraphics[width = 10 cm] {\iftoggle{BW}{shear_completion2-BW}{shear_completion}}
\nota{Near an ideal vertex $v$, horocycles (in \iftoggle{BW}{black}{blue}) and edges of the triangulation (in \iftoggle{BW}{grey}{red}) are mutually orthogonal. They behave much differently in presence of a cusp (left) or of a geodesic boundary (right): in a cusp, horocycles form parallel circles and edges point toward the cusp; near a geodesic boundary, the horocycle point orthogonally toward the boundary and the edges spin and wind around it indefinitely.}
\label{shear_completion:fig}
\end{center}
\end{figure}

If $d_v\neq 0$, the edges of the triangulation pointing towards $v$ spin around the geodesic boundary as shown in Figure \ref{shear_completion:fig}-(right). The spinning direction induces an orientation on each component of $\partial \overline \Sigma$, which depends on the sign of $d_v$. 

Summing up, every shear coordinate $d = (d_1,\ldots, d_{3k})\in \matR^{3k}$ determines a hyperbolic surface $\overline{\Sigma}$ with some $p'$ cusps and $b$ oriented geodesic boundary, where $b+p'=p$ is the number of vertices of the ideal triangulation. The surface $\overline \Sigma$ is diffeomorphic to $S_{g,b,p'}$. Let $\Teich^{\rm or}(S_{g,b,p'})$ denote the Teichm\"uller space of the surface $S_{g,b,p'}$ equipped with an orientation of the $b$ boundary components: of course this is just $2^b$ copies of $\Teich(S_{g,b,p'})$.

\begin{prop} \label{straighten:incomplete:prop}
The shear coordinates induce a bijection
$$\matR^{3k} \longleftrightarrow \bigsqcup_{b+p' = p} \Teich^{\rm or}(S_{g,b,p'}).$$
\end{prop}
\begin{proof}
The map sends $d\in\matR^{3k}$ to the complete hyperbolic $\overline \Sigma$, and we now exhibit an inverse. Let $\overline\Sigma$ be a surface of finite type with a hyperbolic structure with $p'$ cusps and $b$ oriented geodesic boundary components, and $\Sigma = \interior{\overline\Sigma}$. We may straighten any ideal triangulation $\calT$ of $\Sigma$ similarly as we did in Proposition \ref{straightens:prop}, with the only difference that horocycles are replaced by \emph{oriented} geodesic lines as lifts of oriented geodesic boundary components; in the construction, we use the final endpoint of the oriented line instead of the ideal point of the horocycle. We get an ideal triangulation of $\Sigma$ whose completion is $\overline \Sigma$.
\end{proof}

For instance, $\matR^3$ parametrizes altogether all the hyperbolic metrics on the pair of pants, where each boundary component becomes either an oriented geodesic or ``degenerates'' to a cusp.

\subsection{References}
Most of the material presented here is standard and can be found in many books, starting from Thurston's notes \cite{Th}. The main reference is Farb -- Margalit \cite{FM}, which also contains the proof of Proposition \ref{strettamente:convessa:prop}, that appeared originally in a paper of Bestvina, Bromberg, Fujiwara, and Souto \cite{BBFS}. The proof of the Collar Lemma is taken from Hubbard \cite{Hu}.

